\chapter{The config file}\label{ch:config}

ATE reads one machine-wide config file at launch: \texttt{\%PROGRAMDATA\%\textbackslash
ATE\textbackslash config.json} on Windows, \texttt{/etc/ate/config.json} on Linux, or whatever
\texttt{--config PATH} names. It is read \emph{once}. Changing it does nothing until ATE is
restarted; rescanning Test Definitions does not re-read it, and there is no per-user override.

\section{Every key}

Keys are camelCase and grouped in sections. A key that is missing is filled with its default; a key
that is present is checked for type.

\begin{center}
\footnotesize
\begin{tabular}{@{}>{\raggedright\arraybackslash}p{4.4cm}>{\raggedright\arraybackslash}p{3.0cm}>{\raggedright\arraybackslash}p{8.0cm}@{}}
\toprule
Key & Default & Meaning \\
\midrule
\texttt{format} & required & Must be exactly \texttt{1}. A missing or other \texttt{format} stops
  startup. \\
\texttt{defaultProject} & \texttt{""} & Prefills the Project field on the Run screen. A Test Run
  cannot start with an empty Project. \\
\texttt{station.name} & hostname & The Station's descriptive name, recorded on every Test Run. May
  be renamed. \\
\texttt{station.uuid} & generated & The Station's stable identity. If it is missing or empty, ATE
  generates one and writes it back into the file, keeping the other keys and their order. Never copy
  one Station's UUID to another: ATE does not detect a cloned one. \\
\texttt{testDefinitions.}\allowbreak\texttt{directories} & \texttt{<data>/}\allowbreak\texttt{definitions} & Directories searched
  recursively for \texttt{.atedef} files. One that does not exist is a warning, not a failure. \\
\texttt{database.type} & \texttt{sqlite} & v1 accepts only \texttt{sqlite}; anything else stops
  startup. \\
\texttt{database.path} & \texttt{<data>/ate.db} & The SQLite file. Created and upgraded by
  \texttt{ate db migrate} (chapter~\ref{ch:database}). \\
\texttt{database.}\allowbreak\texttt{storeMismatchedFiles} & \texttt{ask} & \texttt{ask}, \texttt{always} or
  \texttt{never}: whether the raw bytes of a Lua file whose checksum no longer matches are stored
  with the Test Run (chapter~\ref{ch:media-releases}). \\
\texttt{reports.directory} & \texttt{<data>/reports} & Where Test Reports are written. \\
\texttt{reports.organisation} & \texttt{""} & Text in every Test Report page header. \\
\texttt{reports.logo} & \texttt{null} & The validator accepts only \texttt{null} today: a path is
  refused with \texttt{must be a null}, even though the report renderer can embed one. Setting a
  logo is not currently possible through the config. \\
\texttt{reports.paperSize} & \texttt{A4} & \texttt{A4} or \texttt{Letter}. The config does not check
  the value; the PDF helper refuses anything else when a report is rendered. \\
\texttt{log.directory} & \texttt{<data>/logs} & Where the application log is written
  (chapter~\ref{ch:logs}). \\
\texttt{log.level} & \texttt{info} & Floor for the application log: \texttt{debug}, \texttt{info},
  \texttt{warning} or \texttt{error}. Anything else is treated as \texttt{info}. \\
\texttt{devices.}\allowbreak\texttt{receiveBufferBytes} & 16777216 & Per-Device receive buffer cap, in bytes. Must be
  positive; overflow drops the oldest bytes. \\
\texttt{deviceTrace.enabled} & \texttt{false} & One \texttt{DEBUG} log line per Device operation, in
  the Run log. \\
\texttt{deviceTrace.maxBytes} & 256 & Traced bytes are shortened past this. \\
\texttt{ui.theme} & \texttt{system} & Accepted and stored; nothing in the shipped code reads it. \\
\texttt{ui.layout} & \texttt{staged} & The Run screen layout: \texttt{staged} or
  \texttt{threePane}. Anything else is refused as a typo. \\
\texttt{ui.allowLayoutSwitch} & \texttt{true} & Whether the layout can be changed for the session. \\
\texttt{ui.allowHistoryDuringRun} & \texttt{true} & Whether History can be browsed while a Test Run
  is in progress. \\
\bottomrule
\end{tabular}
\end{center}

\section{Paths in the config}

A path is absolute or relative; a relative one is resolved against the config file's own directory.
\texttt{<data>} stands for the platform data directory --- \texttt{\%PROGRAMDATA\%\textbackslash
ATE} on Windows, \texttt{/var/lib/ate} on Linux --- and it is the only placeholder that exists. Any
other angle-bracketed token is a config error naming the file and the key, so a typo never becomes
a directory literally called \texttt{<whatever>}. There is no environment-variable expansion and no
\texttt{\textasciitilde} expansion in config \emph{values}; forward slashes work on Windows because
paths are never translated. Only the platform's own default locations expand \texttt{\%PROGRAMDATA\%}.

\section{What makes ATE refuse to start}

An unknown top-level key, a section that is not an object, a value of the wrong type, a bad
\texttt{format}, a \texttt{database.type} that is not \texttt{sqlite}, a
\texttt{storeMismatchedFiles} outside \texttt{ask}/\texttt{always}/\texttt{never}, a
\texttt{ui.layout} that is neither shipped layout, and a \texttt{receiveBufferBytes} of zero or
less all stop startup with a message naming the file and the key. A file that is not valid JSON,
or is not an object at all, is refused the same way.

When the config file is missing, ATE shows the path it expected and asks: create a default config
and launch, or abort. If it cannot write the file --- an Operator without admin rights, say ---
startup is refused with the path. If a config exists but has no \texttt{station.uuid} and the write
back fails, ATE prints the generated UUID for an admin to add by hand and stops.

Two environment variables move the platform defaults, \texttt{ATE\_}\allowbreak\texttt{DEFAULT\_CONFIG}
and \texttt{ATE\_}\allowbreak\texttt{DATA\_DIR}. They exist so the test suite can exercise
"no config at the platform location" without writing to \texttt{/etc/ate}; nothing else in ATE
reads them. Do not use them on a Station.

\cando{write a config file for a Station by hand, and explain exactly what ATE will and will not
accept in it.}
